% !TeX root = ../教材题库_全部条目_转写工作稿.tex
% Source: XJ-B1-C04-RV-P152-Y1473
% 第4章 / 章末复习 复习题四 / 教材 p.152、153 / PDF 第159、160页
% OCR未校对：公式与文字均可能有误；未完成数学TeX转写；禁止直接发给学生
% 原题及相邻原解截图（核对用）：../题目原图/XJ-B1-C04-RV-P152-Y1473-01.webp, ../题目原图/XJ-B1-C04-RV-P152-Y1473-02.webp
% 本文件只在转写完成后作为题干使用；OCR 条目尚不是可靠题干。
\begin{bankitem}{复习题四 第18题}{第4章 / 章末复习 复习题四 / 教材 p.152、153 / PDF 第159、160页}{XJ-B1-C04-RV-P152-Y1473}{OCR未校对：公式与文字均可能有误；未完成数学TeX转写；禁止直接发给学生}
\begin{ocrdraft}
[教材p.152 / PDF第159页]
\par
18．已知幂函数不一（m0）．
\par
（1）若不x）的图象在xe（0，I)时位于直线3，一的上方，求实数m的取值
\par
汜围；
\par
[教材p.153 / PDF第160页]
\par
（2）
\par
若丆）的图象在（1，十。。）时位于直线3，一x的上方，
\par
求实数m的取
\par
值范围．
\par
3X
\end{ocrdraft}
\end{bankitem}
